\section{Lecture 23 (March 10, 2025, est.): Beyond fusion categories}
\label{lec23}
\noindent\textit{Scribe: Guillermo Sanmarco}

This final lecture describes the extension of the Frobenius method to
nonsemisimple categories, and the role of higher Verlinde categories.
Throughout, $k$ is algebraically closed of characteristic $p>0$, and
$\mathcal C$ is a symmetric tensor category with finite-dimensional
Hom-spaces and finite-length objects. Recall that \emph{moderate growth}
means: for every $X$ there is a real constant $d_X$ such that
$\ell(X^{\otimes n})\leq d_X^n$ for every $n\geq1$.
The results and conjectures below retain the distinction made in the
course between proved structure theorems and the proposed higher theory.

\subsection{Frobenius functors in the general case}
For $G\leq S_n$, the symmetry gives a functor
\[
 P_G:\mathcal C\longrightarrow\Rep(G,\mathcal C),\qquad
 X\longmapsto X^{\otimes n}
\]
with the restricted permutation action. There is a general construction
\cite[Section 3.1]{CEOFrobenius}
\begin{equation}\label{lec23-general-frobenius}
 \mathrm{Fr}_G:\mathcal C\longrightarrow
 \mathcal C\boxtimes\overline{\Rep_k(G)}.
\end{equation}
Its second step applies semisimplification in the $G$-representation
variable. When $\mathcal C$ is not semisimple, this needs care:
$\Rep(G,\mathcal C)$ is not generally
$\mathcal C\boxtimes\Rep_k(G)$, and quotienting \emph{all} negligible
morphisms of $\Rep(G,\mathcal C)$ would also lose information from
$\mathcal C$.

One precise construction presents $\mathcal C$ as comodules of a
cotriangular copseudo-Hopf algebra $H$. Such a presentation records the
abelian category by a coalgebra and its tensor structure by compatible
bicomodule data. Then
\[
 \Rep(G,\mathcal C)\simeq
 H\text{-}\mathrm{Comod}(\Rep_k(G)).
\]
Apply $\Rep_k(G)\to\overline{\Rep_k(G)}$ to the representation-variable
data and the coactions. The result is
$H\text{-}\mathrm{Comod}(\overline{\Rep_k(G)})
\simeq\mathcal C\boxtimes\overline{\Rep_k(G)}$.
This construction is independent of the chosen presentation, and is
symmetric monoidal; these facts are part of the cited construction.

The functor \eqref{lec23-general-frobenius} is additive precisely when a
Sylow $p$-subgroup of $G$ acts transitively on $\{1,\ldots,n\}$.
In particular, additivity requires $n=p^j$. In this case it is
$p^j$-semilinear and becomes $k$-linear after replacing the first factor
of its target by $\mathcal C^{(j)}$, with scalar action
$\lambda\cdot f=\lambda^{p^j}f$.

The principal examples are
\[
 \mathrm{Fr}=\mathrm{Fr}_{C_p}:
 \mathcal C\longrightarrow\mathcal C^{(1)}\boxtimes\mathsf{Ver}_p
\]
and the \emph{enhanced Frobenius functor}
\[
 \mathrm{Fr}^{\mathrm{en}}=\mathrm{Fr}_{S_p}:
 \mathcal C\longrightarrow
 \mathcal C^{(1)}\boxtimes\overline{\Rep_k(S_p)}.
\]
Neither is assumed exact in a general category.

\subsection{Invariants, coinvariants, and the functor $\mathrm{Fr}_+$}
For $m\geq1$, write
\[
 \Gamma^mX=(X^{\otimes m})^{S_m},\qquad
 S^mX=(X^{\otimes m})_{S_m}.
\]
These are respectively the largest subobject on which all permutations
act trivially and the largest quotient with trivial permutation action.
The inclusion followed by the quotient gives a canonical map
$\Gamma^mX\to S^mX$.

\begin{definition}
For $j\geq1$, define
\[
 \mathrm{Fr}^{(j)}_+(X)=
 \operatorname{im}\bigl(\Gamma^{p^j}X\longrightarrow S^{p^j}X\bigr),
 \qquad \mathrm{Fr}_+=\mathrm{Fr}^{(1)}_+.
\]
\end{definition}
These are $p^j$-semilinear functors. The functor $\mathrm{Fr}_+$ is a
direct summand of $\mathrm{Fr}^{\mathrm{en}}$, extracted by the unit
multiplicity in $\overline{\Rep_k(S_p)}$. It has a canonical lax monoidal
structure; its structural maps are not asserted to be isomorphisms.

\begin{example}
For an ordinary vector space $V$ with basis $v_1,\ldots,v_d$, the image of
$\Gamma^{p^j}V\to S^{p^j}V$ has basis $v_1^{p^j},\ldots,v_d^{p^j}$.
Indeed, the orbit sum of a mixed tensor maps to its orbit size times the
corresponding monomial. The multinomial coefficient for a nonconstant
word of length $p^j$ is divisible by $p$, whereas constant words have
orbit size one. Thus this functor recovers the usual Frobenius twist of
a vector space, and similarly the Frobenius twist of a representation.
\end{example}

\begin{theorem}[Coulembier--Etingof--Ostrik]
\label{lec23-frobenius-exact}
The following conditions are equivalent:
\begin{enumerate}
 \item $\mathrm{Fr}$ is exact;
 \item $\mathrm{Fr}^{\mathrm{en}}$ is exact;
 \item $\mathrm{Fr}_+$ is exact.
\end{enumerate}
A category satisfying them is called \emph{Frobenius exact}.
\end{theorem}
See \cite[Proposition 3.4]{CEOFrobenius}. The equivalence is a theorem,
not a formal consequence of additivity. In particular,
semisimple categories are Frobenius exact, but nonsemisimple examples
need not be.

\begin{theorem}[The Verlinde fiber-functor criterion]
\label{lec23-ceo-fiber}
The category $\mathcal C$ is Frobenius exact and of moderate growth if
and only if it admits an exact $k$-linear symmetric monoidal functor
$\mathcal C\to\mathsf{Ver}_p$.
\end{theorem}
This is the general structure theorem of
\cite[Theorem 7.13]{CEOFrobenius}. The forward implication is a deep
extension of the fusion-category argument. For the reverse implication,
exact faithful functors preserve a strict composition series and hence
cannot decrease length. Tensor powers in the fusion category
$\mathsf{Ver}_p$ have exponentially bounded length, proving moderate
growth. Naturality of the Frobenius construction and exact faithfulness
of the fiber functor then detect exactness of $\mathrm{Fr}$, since the
target is semisimple.

\begin{lemma}[Iteration]\label{lec23-iteration}
If $\mathcal C$ is Frobenius exact, then
\[
 \mathrm{Fr}^{(j)}_+\simeq(\mathrm{Fr}_+)^{\circ j}
 \qquad(j\geq1).
\]
\end{lemma}
This is \cite[Lemma 5.1]{CEOFrobenius}; the equality compares the same
semilinearity on both sides. Its proof uses a faithful exact symmetric
functor that splits all short exact sequences, and the comparison of
invariants and coinvariants for iterated permutation groups. Without
Frobenius exactness the assertion can fail. Thus the definition using
$S_{p^j}$-invariants must not be replaced by iteration without this
hypothesis.

\subsection{Higher Verlinde categories and incompressibility}
\begin{definition}
A symmetric tensor category $\mathcal D$ is \emph{incompressible} if
every tensor functor from $\mathcal D$ to another symmetric tensor
category is an embedding of a tensor subcategory: it is fully faithful
and its image is closed under subobjects and quotients.
\end{definition}
Besides $\Vect$, $\SVec$ for $p>2$, and $\mathsf{Ver}_p$, there are
finite symmetric tensor categories $\mathsf{Ver}_{p^n}$ and
$\mathsf{Ver}_{p^n}^+$ for $n\geq1$, constructed from quotients of the
category of tilting modules for $\SL_2$ and their abelian envelopes
\cite{BEOHigher}. There are tensor embeddings
\[
 \mathsf{Ver}_p\subseteq\mathsf{Ver}_{p^2}
 \subseteq\mathsf{Ver}_{p^3}\subseteq\cdots.
\]
These categories and their indicated subcategories are incompressible.
For $n>1$, $\mathsf{Ver}_{p^n}$ is not Frobenius exact. In particular,
moderate growth alone does not imply the existence of a fiber functor
to $\mathsf{Ver}_p$.

\begin{theorem}[Coulembier--Etingof--Ostrik]
Every symmetric tensor category of moderate growth admits a tensor
functor to an incompressible symmetric tensor category.
\end{theorem}
This existence theorem is proved in \cite{CEOIncompressible}. It leaves
the classification of possible incompressible targets as a separate
problem. The conjecture discussed in the course is that every category
of moderate growth fibers over
$\mathsf{Ver}_{p^\infty}=\bigcup_n\mathsf{Ver}_{p^n}$; the known higher
Verlinde categories and their tensor subcategories would then supply
all the incompressible targets in this class. This is stated as a
conjecture, not used as a theorem.

\subsection{A proposed higher Frobenius construction}
We conclude with the construction of Coulembier--Flake
\cite{CFHigher}. Fix $n\geq1$ and $E=(C_p)^n$. Choose a faithful
two-dimensional representation $V$ of $E$, equivalently an embedding
$E\hookrightarrow\SL_2(k)$. For example, choose
$\mathbb F_p$-linearly independent elements $a_1,\ldots,a_n\in k$ and
send the $i$th generator to
$\left(\begin{smallmatrix}1&a_i\\0&1\end{smallmatrix}\right)$.
Let $\mathcal D$ be the full additive idempotent-complete monoidal
subcategory generated by $V$ in $\Rep_k(E)$.

There is a tensor ideal of morphisms $\mathcal K$ in $\mathcal D$ such
that $\mathsf{Ver}_{p^n}$ is the monoidal abelian envelope of
$\mathcal D/\mathcal K$. This is an actual construction theorem.
For $n=1$ it recovers the quotient of $\Rep_k(C_p)$ by negligible maps.
Together with the regular permutation embedding $E\hookrightarrow S_{p^n}$,
it leads to candidate functors
\[
 \Phi^V_{\mathcal C}:\mathcal C\longrightarrow
 \mathcal C^{(n)}\boxtimes\mathsf{Ver}_{p^n},
\]
based on the $E$-action on $X^{\otimes p^n}$. The exponent in the
Frobenius scalar twist is $p^n$.

The category $\mathcal D$ has only finitely many indecomposables and has
a smallest nonzero \emph{thick tensor ideal of objects} $\mathcal J$.
Here ``thick'' means closed under finite direct sums and direct summands;
being a tensor ideal means closure under tensoring with every object of
the ambient category. This ideal of objects is distinct from the ideal
of morphisms $\mathcal K$ above.

The two conjectural steps discussed in the handwritten notes are:
\begin{enumerate}
 \item The candidate functors $\Phi^V_{\mathcal C}$ admit compatible
 symmetric monoidal structures and intertwine tensor functors in
 $\mathcal C$, giving higher analogues of the ordinary Frobenius functor.
 \item The full subcategory $\mathcal J\subseteq\mathcal D\subseteq
 \Rep_k(E)$ remains a thick tensor ideal in the entire category
 $\Rep_k(E)$, including closure under tensoring with modules outside
 $\mathcal D$.
\end{enumerate}
Coulembier--Flake prove that the second statement implies the first.
They establish special cases, but the general monoidality assertion is
not a consequence merely of taking a tensor power. The longer-term aim
is an exactness criterion for fibering over $\mathsf{Ver}_{p^n}$,
parallel to Theorem~\ref{lec23-ceo-fiber} when $n=1$.
